\chapter{Multiple Integrals \& Field Theorems}
\index{Multiple Integrals}
\index{Jacobian Determinant}
\index{Line Integrals}
\index{Surface Integrals}
\index{Green's Theorem}
\index{Stokes' Theorem}
\index{Gauss's Divergence Theorem}
\index{Maxwell's Equations}

\begin{chapterobjectives}
Upon mastering this chapter, students will be able to:
\begin{enumerate}
    \item Evaluate multiple integrals over general domains using Fubini's theorem.
    \item Perform coordinate transformations in multivariable integrals using the Jacobian determinant.
    \item Calculate line integrals of vector fields along specified paths and surface fluxes through arbitrary surfaces.
    \item State, prove, and apply Green's theorem, Stokes' circulation theorem, and Gauss's Divergence theorem.
    \item Connect differential and integral formulations of physical laws (Maxwell's equations, continuity equation).
    \item Implement numerical multiple integration algorithms in Python.
\end{enumerate}
\end{chapterobjectives}

\section{Multiple Integrals and the Jacobian Transformation}
Physical quantities such as total mass $M = \iiint \rho(\vec{r})\,dV$, center of mass $\vec{r}_{\text{cm}} = \frac{1}{M}\iiint \vec{r}\rho\,dV$, and moments of inertia require integration over two- and three-dimensional spatial manifolds.

When changing integration variables from $(x, y, z)$ to curvilinear coordinates $(u, v, w)$, the volume transformation is scaled by the absolute value of the \textbf{Jacobian determinant}:
\begin{equation}
J \equiv \frac{\partial(x, y, z)}{\partial(u, v, w)} = \begin{vmatrix} \frac{\partial x}{\partial u} & \frac{\partial x}{\partial v} & \frac{\partial x}{\partial w} \\ \frac{\partial y}{\partial u} & \frac{\partial y}{\partial v} & \frac{\partial y}{\partial w} \\ \frac{\partial z}{\partial u} & \frac{\partial z}{\partial v} & \frac{\partial z}{\partial w} \end{vmatrix}
\end{equation}
The integral transformation law is:
\begin{equation}
\iiint_V f(x, y, z)\,dx\,dy\,dz = \iiint_{V'} f(x(u,v,w), y(u,v,w), z(u,v,w))\,|J|\,du\,dv\,dw
\end{equation}
For orthogonal curvilinear coordinates, the Jacobian equals the product of the scale factors: $|J| = h_1 h_2 h_3$. Specifically:
\begin{itemize}
    \item Cylindrical: $|J| = \rho \implies dV = \rho\,d\rho\,d\phi\,dz$.
    \item Spherical: $|J| = r^2\sin\theta \implies dV = r^2\sin\theta\,dr\,d\theta\,d\phi$.
\end{itemize}

\section{The Fundamental Theorems of Vector Integral Calculus}

\begin{maththeorem}[Green's Theorem in the Plane]
Let $C$ be a positively oriented (counterclockwise), piecewise-smooth, simple closed curve enclosing region $D \subset \mathbb{R}^2$. If $P(x, y)$ and $Q(x, y)$ have continuous partial derivatives on $D$:
\begin{equation}
\oint_C (P\,dx + Q\,dy) = \iint_D \left(\frac{\partial Q}{\partial x} - \frac{\partial P}{\partial y}\right)dA
\end{equation}
\end{maththeorem}

\begin{maththeorem}[Stokes' Theorem (Circulation Theorem)]
Let $S$ be an oriented piecewise-smooth surface bounded by a closed contour $C = \partial S$ oriented by the right-hand rule with respect to the unit normal $\hat{n}$. For any continuously differentiable vector field $\vec{A}$:
\begin{equation}
\oint_C \vec{A} \cdot d\vec{r} = \iint_S (\nabla \times \vec{A}) \cdot d\vec{a}
\end{equation}
\end{maththeorem}
In electrodynamics, applying Stokes' theorem to Faraday's law $\nabla \times \vec{E} = -\frac{\partial \vec{B}}{\partial t}$ generates the macroscopic induction law:
\begin{equation}
\mathcal{E} = \oint_C \vec{E} \cdot d\vec{r} = -\frac{d}{dt}\iint_S \vec{B} \cdot d\vec{a} = -\frac{d\Phi_B}{dt}
\end{equation}

\begin{maththeorem}[Gauss's Divergence Theorem]
Let $V$ be a compact three-dimensional region bounded by a closed surface $S = \partial V$ oriented with an outward-pointing unit normal $\hat{n}$. For any continuously differentiable vector field $\vec{A}$:
\begin{equation}
\oiint_S \vec{A} \cdot d\vec{a} = \iiint_V (\nabla \cdot \vec{A})\,dV
\end{equation}
\end{maththeorem}
Gauss's divergence theorem bridges the microscopic differential form of Gauss's law $\nabla \cdot \vec{E} = \frac{\rho}{\epsilon_0}$ and its macroscopic integral form $\oiint_S \vec{E} \cdot d\vec{a} = \frac{Q_{\text{enclosed}}}{\epsilon_0}$.

\begin{pitfallbox}[Orientation and Right-Hand Rule Conventions]
In Stokes' theorem, the boundary curve $C$ must be traversed such that when walking along $C$ with one's head pointing in the direction of the surface normal $\hat{n}$, the surface $S$ lies to one's left. In Gauss's theorem, $S$ must be a \textbf{closed} surface, and $d\vec{a} = \hat{n}\,da$ must point \textbf{outward}.
\end{pitfallbox}

\section{Worked Examples and Physical Applications}

\begin{psctexample}[Verification of Gauss's Divergence Theorem]
Let $\vec{F}(x, y, z) = x\hat{i} + y\hat{j} + z\hat{k} = \vec{r}$. Verify Gauss's divergence theorem over a solid sphere of radius $R$ centered at the origin.

\textbf{\color{rbruNavy}Picture / Conceptualization:}
The vector field points radially outward everywhere from the origin with magnitude proportional to distance. The bounding surface is the sphere $S$ of radius $R$.

\textbf{\color{rbruNavy}Strategy / Governing Laws:}
Evaluate both sides of $\oiint_S \vec{F} \cdot d\vec{a} = \iiint_V (\nabla \cdot \vec{F})\,dV$ independently and compare the results.

\textbf{\color{rbruNavy}Calculation / Analytical Solution:}
1. Volume integral:
Compute the divergence of the position vector:
\begin{equation}
\nabla \cdot \vec{r} = \frac{\partial x}{\partial x} + \frac{\partial y}{\partial y} + \frac{\partial z}{\partial z} = 1 + 1 + 1 = 3
\end{equation}
Integrate over the spherical volume:
\begin{equation}
\iiint_V (\nabla \cdot \vec{F})\,dV = 3 \iiint_V dV = 3\left(\frac{4}{3}\pi R^3\right) = 4\pi R^3
\end{equation}

2. Surface integral:
On the spherical surface of radius $R$, the outward unit normal is $\hat{n} = \hat{r}$, and the vector field is $\vec{F} = R\hat{r}$.
The surface area element is $d\vec{a} = \hat{r}\,da$.
\begin{equation}
\oiint_S \vec{F} \cdot d\vec{a} = \oiint_S (R\hat{r}) \cdot (\hat{r}\,da) = R \oiint_S da = R (4\pi R^2) = 4\pi R^3
\end{equation}

\textbf{\color{rbruNavy}Think / Physical Insights:}
Both calculations yield precisely $4\pi R^3$. The divergence theorem converted a two-dimensional spherical surface integral into a trivial volume integral of a constant scalar divergence!
\qedexam
\end{psctexample}

\begin{pythonbox}[Numerical Verification of Flux through a Sphere]
import numpy as np

# Monte Carlo integration over unit sphere
N = 1000000
pts = np.random.uniform(-1, 1, (N, 3))
r_sq = np.sum(pts**2, axis=1)
inside = pts[r_sq <= 1.0]

vol_sphere = (len(inside) / N) * 8.0 # Box volume is 2^3 = 8
div_F = 3.0
vol_integral = div_F * vol_sphere

exact = 4.0 * np.pi * (1.0)**3
print(f"Numerical Volume Integral: {vol_integral:.4f}")
print(f"Exact Analytical Flux:     {exact:.4f}")
\end{pythonbox}

\section{Summary of Key Formulas}
\begin{table}[H]
\centering
\small
\begin{tabularx}{\textwidth}{llX}
\toprule
\textbf{Theorem} & \textbf{Integral Equation} & \textbf{Physical Meaning} \\
\midrule
Green's Theorem & $\oint_{\partial D} (P dx + Q dy) = \iint_D (\frac{\partial Q}{\partial x} - \frac{\partial P}{\partial y}) dA$ & Planar circulation equals enclosed vorticity flux \\
Stokes' Theorem & $\oint_{\partial S} \vec{A} \cdot d\vec{r} = \iint_S (\nabla \times \vec{A}) \cdot d\vec{a}$ & Faraday's induction and Ampère's circuital laws \\
Divergence Theorem & $\oiint_{\partial V} \vec{A} \cdot d\vec{a} = \iiint_V (\nabla \cdot \vec{A}) dV$ & Gauss's law for electrostatic and gravitational flux \\
Jacobian Rule & $dV = |\det J|\,du_1 du_2 du_3$ & Geometric volume distortion under curvilinear transforms \\
\bottomrule
\end{tabularx}
\caption{Summary of fundamental integral transformation theorems.}
\label{tab:ch07_summary}
\end{table}

\section{Chapter Exercises}
\begin{enumerate}
    \item \textbf{Stokes' Theorem on a Paraboloid:} Verify Stokes' theorem for $\vec{A} = -y\hat{i} + x\hat{j} + z\hat{k}$ over the paraboloid surface $z = 1 - x^2 - y^2$ ($z \ge 0$) bounded by the unit circle in the $xy$-plane.
    \item \textbf{Gravitational Field Divergence:} Show that for the Newtonian gravitational field $\vec{g} = -\frac{GM}{r^2}\hat{r}$, $\oiint_S \vec{g} \cdot d\vec{a} = -4\pi GM$ for any closed surface enclosing the origin.
    \item \textbf{Jacobian in Elliptic Coordinates:} Calculate the Jacobian determinant for planar elliptic coordinates $x = a\cosh u\cos v$, $y = a\sinh u\sin v$.
\end{enumerate}
